\section{Discussion}
\label{sec:discussion}

\subsection{What this design can and cannot identify}
\label{sec:identifiability}

The question that motivated this work --- how much of a web agent's robustness is
attributable to the model versus to the architecture --- cannot be answered by
these data, and we want to be precise about why rather than gesturing at
``small $n$.''

Identifying a model effect requires the model to vary. In the official subset and
the behavioural matrix the model is constant; in the architecture pilot it is also
constant. A variance decomposition over these factors would partition variance
into a term that is zero by construction, and any numeric ``share'' it produced
would be an artefact of which factor happened to be varied. The same argument
applies to the model$\times$architecture interaction.

What \emph{is} identified here is the fault effect, within one
(model, architecture) configuration, and --- more usefully --- the properties of
the measurement. Both of those are estimated with a paired design and reported
with the power to interpret them.

The design that would identify the attribution question is a crossed one:
$\mathrm{Model} \times \mathrm{Architecture} \times \mathrm{Fault}$, with the same
paired control/fault structure inside each cell, evaluated natively. Two models
and two architectures over the five faults on the same \OfficialTasks{} tasks
would be enough to estimate main effects and the interaction; more models would be
needed to separate a model effect from an idiosyncratic one. We give the concrete
staged plan, with the measured per-trial wall-clock that sizes it, in the
supplementary experiment plan released with this paper.

\subsection{Why the nulls are the finding}

It would be easy to read \secref{sec:res-main} as a failed experiment and either
bury it or explain it away. We argue the opposite: the nulls are reproducible and
they have a structural cause. The success outcome is binary, the tasks are
heterogeneous, the per-fault discordance is only \PiDNative{}, and the plausible
effect sizes (a few percentage points) are far below what the design can resolve.
Any underpowered fault-injection study on WebArena will produce this pattern, and
because the pattern admits a flattering reading --- ``no significant degradation,
therefore robust'' --- it is likely to be published as a positive result. Naming
the failure mode seems more useful than adding one more underpowered table.

The constructive version of the same point is that the process measures
\emph{do} have power here. Step deltas of $2$--$4$ are large relative to their
within-pair variance, and they are significant at $n=\BehWebHttpErrorN{}$ pairs
where the binary outcome is not. If robustness is the object of study, the
measurement should be chosen with that in mind.

\subsection{Threats to validity}
\label{sec:threats}

\paragraph{Instrumentation.} The dominant threat is the evaluator fallback
(\secref{sec:fallback}), which we handle by stratification rather than by
correction. A second risk is that the native path may itself mis-score tasks in
ways we have not audited; the behavioural matrix, which has no official verdict at
all, is reported separately for exactly this reason.

\paragraph{Synthetic faults.} Our faults are injected at a fixed step by a seeded
injector: deterministic and reproducible, which is what makes the paired analysis
possible, but not draws from a real fault distribution. The injection \emph{step}
in particular is a free parameter we held fixed, and an agent's robustness
plausibly depends on when the fault arrives. This is the confound a reviewer
should press on hardest, and the first thing our staged plan adds.

\paragraph{Scope.} The behavioural matrix uses \MatrixTasks{} tasks and
\TrialsPerCell{} seeds --- a larger and different task set than the
\OfficialTasks{}-task official subset, and not joinable to it. The process results
therefore generalise over more tasks than the success results do. As shown in
\secref{sec:res-recovery}, the recovery artifact is degenerate and we draw no
recovery conclusions from it; a credible analysis would need newly defined,
outcome-independent labels, which is a contribution in its own right. Finally,
one model, one architecture, one temperature, one provider: our claims are about
the measurement protocol and this configuration's fault sensitivity, not about web
agents in general.
